\documentclass[a4paper, 12pt]{article}
\usepackage[T2A]{fontenc}
\usepackage[english,russian]{babel}
\usepackage{amsmath}
\usepackage{amsthm}
\usepackage{amsfonts}
\usepackage{amssymb}
\usepackage{pgfplots}

\usepackage{enumitem}
% Абзацы и геометрия страницы
\usepackage{indentfirst}
%\usepackage{geometry}
\usepackage[
    top=2cm,
    bottom=2cm,
    left=1cm,
    right=1cm
]{geometry}
\begin{document}


\begin{center}
    МИНИСТЕРСТВО НАУКИ И ВЫСШЕГО ОБРАЗОВАНИЯ РОССИЙСКОЙ ФЕДЕРАЦИИ

    ГОСУДАРСТВЕННОЕ ОБРАЗОВАТЕЛЬНОЕ УЧРЕЖДЕНИЕ ВЫСШЕГО ПРОФЕССИОНАЛЬНОГО ОБРАЗОВАНИЯ

    РЯЗАНСКИЙ ГОСУДАРСТВЕННЫЙ РАДИОТЕХНИЧЕСКИЙ УНИВЕРСИТЕТ

    ИМ. В. Ф. УТКИНА

    КАФЕДРА ВЫЧИСЛИТЕЛЬНОЙ И ПРИКЛАДНОЙ МАТЕМАТИКИ
\end{center}

\vspace{5 cm}

\begin{center}
    ЛАБОРАТОРНАЯ РАБОТА № 6

    ПО ДИСЦИПЛИНЕ "ТЕОРЕТИЧЕСКАЯ ИНФОРМАТИКА"

    НА ТЕМУ "ПРЕДСТАВЛЕНИЕ ЧИСЕЛ С ПЛАВАЮЩЕЙ ТОЧКОЙ"

    ВАРИАНТ 14

\end{center}

\vspace{7 cm}

\begin{flushright}
    Выполнил:

    Студент группы 6413

    Косяков Д.А.

    Проверил

    Доцент кафедры ВПМ

    Никичкин Борис Владимирович
\end{flushright}

\begin{center}
    РЯЗАНЬ 2026
\end{center}

\pagebreak

\textbf{ Цель работы:} определить, какое число представлено типом single (число с плавающей точкой) по его побайтовому шестнадцатеричному представлению.

\textbf{ Задание:} дано побайтовое шестнадцатеричное представление типа single.
Определить, какое число представлено этим представлением.

\textbf{ Выполнение задания:} дано шестнадцатеричное представление $\mathrm{424A0000}_{16}$.

Переведем представление в двоичную систему счисления. Каждую шестнадцатеричную цифру
заменяем четырьмя двоичными разрядами ($4 = 0100$, $2 = 0010$, $A = 1010$, $0 = 0000$):
$$01000010 \quad 01001010 \quad 00000000 \quad 00000000$$

Разделим двоичное представление на знак (1 бит), порядок (8 бит) и мантиссу (23 бита):
$$0 \quad 10000100 \quad 10010100000000000000000$$

Знак равен 0, значит число положительное.

Порядок равен $10000100_2 = 1 \cdot 2^7 + 1 \cdot 2^2 = 128 + 4 = 132$. Скорректируем порядок
(вычтем смещение 127):
$$132 - 127 = 5$$

К мантиссе добавим слева скрытую единицу (показана жирным шрифтом) и перенесем запятую от скрытой
единицы вправо на полученную величину порядка (порядок равен 5):
$$\mathbf{1},10010100000000000000000 \;\rightarrow\; \mathbf{1}10010{,}100000000000000000$$

Целая часть числа равна
$$110010_2 = 1 \cdot 2^5 + 1 \cdot 2^4 + 0 \cdot 2^3 + 0 \cdot 2^2 + 1 \cdot 2^1 + 0 \cdot 2^0 = 32 + 16 + 2 = 50$$

Дробная часть числа равна
$$0{,}1_2 = 1 \cdot 2^{-1} = 0{,}5$$

Сложим целую и дробную части числа, знак плюс:
$$+(50 + 0{,}5) = 50{,}5$$

В результате получено число 50,5, которое представляется типом single, как $424A0000$.

Проверка: $50{,}5 = 1{,}578125 \cdot 2^5$; $1{,}578125 = 1{,}100101_2$; порядок $5 + 127 = 132 = 10000100_2$;
собираем $0\;10000100\;10010100000000000000000 = \mathrm{424A0000}_{16}$. Результат совпадает.
\vspace{1cm}

\textbf{Отчет:}
\begin{enumerate}[nosep]
    \item Исходное представление: $\mathrm{424A0000}_{16}$.
    \item Двоичная запись: 0 10000100 10010100000000000000000.
    \item Знак: 0 (плюс), порядок: $132 - 127 = 5$, мантисса со скрытой единицей: 1,1001010\ldots
    \item Число: $110010{,}1_2 = 50{,}5_{10}$.
\end{enumerate}

\end{document}
